\subsection{\texorpdfstring{The case \(\chi = 1\)}{The case \textbackslash chi = 1}}

If \(f\) is a function on \(\mathbf{X} / \mathbf{H}\) , then \(f \circ \pi\) is a function on \(\mathbf{X}\) constant on the orbits, continuous if and only if \(f\) is continuous. The mapping \(f \mapsto f \circ \pi\) defines in particular a bijection of \(\mathcal{K}(\mathbf{X} / \mathbf{H})\) onto \(\mathcal{K}^1(\mathbf{X})\) .

We can then, in the case that \(\chi = 1\) , reformulate certain results of No.~1 in the following way:

Let \(f\) be a continuous numerical function on \(X\) whose support has compact intersection with the saturation of every compact subset of \(X\) . The formula

\[
f ^ {b} (\pi (x)) = \int_ {\mathrm{H}} f (x \xi) d \beta (\xi)\tag{3}
\]

defines a continuous function \(f^{b}\) on \(\mathbf{X} / \mathbf{H}\) . If \(g\) is a continuous function on \(\mathbf{X} / \mathbf{H}\) , then

\[
(f \cdot g \circ \pi) ^ {b} = f ^ {b} \cdot g.\tag{4}
\]

If \(\eta \in \mathbf{H}\) , then

\[
\left(\delta (\eta) f\right) ^ {b} = \Delta_ {\mathrm{H}} (\eta) ^ {- 1} f ^ {b}.\tag{5}
\]

One must not forget that the definition of \(f^{b}\) depends on the choice of \(\beta\) . If H is compact and \(\beta\) is normalized, the function \(f^{b}\) is sometimes called the orbital mean of f.

If \(f \in \mathcal{K}(\mathrm{X})\) , then \(f^{b} \in \mathcal{K}(\mathrm{X / H})\) . The mapping \(f \mapsto f^{b}\) of \(\mathcal{K}(\mathrm{X})\) into \(\mathcal{K}(\mathrm{X / H})\) is linear, and the image of \(\mathcal{K}(\mathrm{X})\) (resp. \(\mathcal{K}_{+}(\mathrm{X})\) ) is \(\mathcal{K}(\mathrm{X / H})\) (resp. \(\mathcal{K}_{+}(\mathrm{X / H})\) ).

Remark 1. --- We are going to show that the mapping \(f \mapsto f^{b}\) is a strict morphism (GT, III, §2, No.~8) of \(\mathcal{K}(\mathrm{X})\) onto \(\mathcal{K}(\mathrm{X}/\mathrm{H})\) .

\begin{enumerate}
\def\labelenumi{\alph{enumi})}
\item
  The mapping is continuous: it suffices to prove that, for every compact subset K of X, the restriction of \(f \mapsto f^{b}\) to \(\mathcal{K}(X, K)\) is a continuous mapping of \(\mathcal{K}(X, K)\) into \(\mathcal{K}(X / H, \pi(K))\) (TVS, II, §4, No.~4, Prop. 5); since H operates properly in X, the set P of \(\xi \in H\) such that K \(\xi\) intersects K is compact; one concludes from (3) that \(\sup_{x \in K} |f^{b}(\pi(x))| \leqslant \beta(P) \sup_{x \in K} |f(x)|\) , and this proves our assertion.
\item
  Let \(K'\) be a compact subset of X/H. Let us choose a compact subset K of X such that \(\pi(K) = K'\) , and let us show that the restriction of \(f \mapsto f^{b}\) to \(\mathcal{K}(X, K)\) is a strict morphism of \(\mathcal{K}(X, K)\) onto \(\mathcal{K}(X/H, K')\) . It suffices to construct a right inverse for this restriction (GT, III, §6, No.~2, Prop. 3). Now, by Lemma 1 of No.~1 (whose notations we adopt), one obtains such an inverse by composing the following mappings:
\end{enumerate}

\(\alpha)\) the mapping \(f' \mapsto f' \circ \pi\) of \(\mathcal{K}(\mathrm{X}/\mathrm{H}, \mathrm{K}')\) into the set E of functions of \(\mathcal{K}^1(\mathrm{X})\) whose support is contained in KH;

\(\beta)\) the mapping of E into E that, to every \(g \in \mathbf{E}\) , makes correspond the function equal to \(g / u^1\) on KH, and to 0 on X - KH;

\(\gamma)\) the mapping of E into \(\mathcal{K}(\mathbf{X})\) that, to every function \(h \in \mathbf{E}\) , makes correspond \(uh\) .

\begin{enumerate}
\def\labelenumi{\alph{enumi})}
\setcounter{enumi}{2}
\tightlist
\item
  This established, if \(\mathbf{V}\) is a convex neighborhood of 0 in \(\mathcal{K}(\mathbf{X})\) , then \(\mathrm{V} \cap \mathcal{K}(\mathrm{X}, \mathrm{K})\) is a convex neighborhood of 0 in \(\mathcal{K}(\mathrm{X}, \mathrm{K})\) , therefore \(\mathrm{V}^{b} \cap \mathcal{K}(\mathrm{X} / \mathrm{H}, \mathrm{K}')\)
\end{enumerate}

is a convex neighborhood of 0 in \(\mathcal{K}(\mathrm{X / H},\mathrm{K}^{\prime})\) by \(b)\) , therefore \(\mathbf{V}^b\) is a neighborhood of 0 in \(\mathcal{K}(\mathrm{X / H})\) (TVS, II, §4, No.~4). This completes the proof.

PROPOSITION 4. --- a) Let \(\lambda\) be a measure on X/H. There exists one and only one measure \(\lambda^{\sharp}\) on X such that

\[
\int_ {\mathrm{X} / \mathrm{H}} f ^ {b} d \lambda = \int_ {\mathrm{X}} f d \lambda^ {\sharp}\tag{6}
\]

for all \(f \in \mathcal{K}(\mathrm{X})\) . One has \(\pmb{\delta}(\xi)\lambda^{\sharp} = \Delta_{\mathrm{H}}(\xi)\lambda^{\sharp}\) for all \(\xi \in \mathrm{H}\) .

\begin{enumerate}
\def\labelenumi{\alph{enumi})}
\setcounter{enumi}{1}
\tightlist
\item
  Conversely, let \(\mu\) be a measure on X such that \(\delta(\xi)\mu = \Delta_{\mathrm{H}}(\xi)\mu\) for all \(\xi \in \mathrm{H}\) . There exists one and only one measure \(\lambda\) on X/H such that \(\mu = \lambda^{\sharp}\) .
\end{enumerate}

This is a special case of No.~1, Prop. 3.

DEFINITION 1. --- With hypotheses and notations as in Prop. 4, \(\lambda\) is called the quotient of \(\mu\) by \(\beta\) and is denoted \(\frac{\mu}{\beta}\) or \(\mu/\beta\) .

The mapping \(\lambda \mapsto \lambda^{\sharp}\) of \(\mathcal{M}(\mathrm{X / H})\) into \(\mathcal{M}(\mathrm{X})\) is none other than the transpose of the mapping \(f\mapsto f^{b}\) of \(\mathcal{K}(\mathrm{X})\) into \(\mathcal{K}(\mathrm{X / H})\) . Let \(\mathfrak{F}\) be a filter on \(\mathcal{M}(\mathrm{X / H})\) ; to say that \(\lim_{\lambda ,\mathfrak{F}}\lambda^{\sharp}(f) = 0\) for all \(f\in \mathcal{K}(\mathrm{X})\) is equivalent to saying that \(\lim_{\lambda ,\mathfrak{F}}\lambda (f^{\prime}) = 0\) for all \(f^{\prime}\in \mathcal{K}(\mathrm{X / H})\) ; therefore the mapping \(\lambda \mapsto \lambda^{\sharp}\) is, for the vague topologies, an isomorphism of \(\mathcal{M}(\mathrm{X / H})\) onto a linear subspace of \(\mathcal{M}(\mathrm{X})\) . This subspace is vaguely closed, since it is the set of \(\mu \in \mathcal{M}(\mathrm{X})\) such that \(\delta (\xi)\mu = \Delta_{\mathrm{H}}(\xi)\mu\) for all \(\xi \in \mathrm{H}\) . It is clear that the conditions \(\lambda \geqslant 0\) and \(\lambda^{\sharp}\geqslant 0\) are equivalent.

The formula (6) may, by analogy with the usual notation for double integrals, be written

\[
\int_ {\mathrm{X}} f (x) d \lambda^ {\sharp} (x) = \int_ {\mathrm{X/H}} d \lambda (\dot {x}) \int_ {\mathrm{H}} f (x \xi) d \beta (\xi) \quad (\dot {x} = \pi (x)).\tag{7}
\]

This involves an abuse of notation, the integral \(\int_{\mathrm{H}} f(x\xi) d\beta(\xi)\) being regarded as a function of \(\dot{x}\) and not of x; this manner of writing will be used frequently in what follows provided no confusion can arise.
% semantic-fragments: legacy-input-seam
\relax{}
